% Part: first-order-logic
% Chapter: sequent-calculus
% Section: proving-things

\documentclass[../../../include/open-logic-section]{subfiles}

\begin{document}

\iftag{FOL}
      {\olfileid{fol}{seq}{pro}}
      {\olfileid{pl}{seq}{pro}}

\olsection{\usetoken{P}{derivation}-এর উদাহরণ}

\begin{ex}
$!A \land !B \Sequent !A$ সিকোয়েন্টটির একটি $\Log{LK}$-!!{derivation} দাও।

!!{derivation}-টির নিচে কাঙ্ক্ষিত অন্তিম সিকোয়েন্টটি লিখে শুরু করি।
\begin{prooftree}
\AxiomC{}
\UnaryInf$!A\land !B \fCenter !A$
\end{prooftree}
এরপর দেখতে হবে, এই রূপের নিচের সিকোয়েন্ট কোন ধরনের অনুমান থেকে আসতে
পারে। সেটি কোনো গঠনগত বিধি হতে পারে, তবে আগে কোনো যৌক্তিক বিধি খোঁজা
ভালো। নিচের সিকোয়েন্টে একমাত্র যৌক্তিক সংযোজক হলো~$\land$। তাই আমাদের
একটি $\land$-বিধি চাই; আর $\land$ চিহ্নটি পূর্বাংশে থাকায় বিধিটি হবে
\LeftR{\land}।
\begin{prooftree}
\AxiomC{}
\RightLabel{\LeftR{\land}}
\UnaryInf$!A\land !B \fCenter !A$
\end{prooftree}
\LeftR{\land} অনুমানের উপরের সিকোয়েন্টের জন্য দুটি সম্ভাবনা আছে:
$!A \Sequent !A$ অথবা $!B \Sequent !A$। স্পষ্টতই $!A \Sequent !A$
একটি প্রারম্ভিক সিকোয়েন্ট (এটি সুবিধাজনক), কিন্তু সাধারণভাবে
$!B \Sequent !A$ নিষ্পাদনযোগ্য নয়। এবার উপরের সিকোয়েন্টটি বসাই:
\begin{prooftree}
\Axiom$!A \fCenter !A$
\RightLabel{\LeftR{\land}}
\UnaryInf$!A\land !B \fCenter !A$
\end{prooftree}
এখন $!A \land !B \Sequent !A$ সিকোয়েন্টটির একটি সঠিক
$\Log{LK}$-!!{derivation} পাওয়া গেল।
\end{ex}

\begin{ex}
$\lnot !A \lor !B \Sequent !A \lif !B$ সিকোয়েন্টটির একটি
$\Log{LK}$-!!{derivation} দাও।

!!{derivation}-টির নিচে কাঙ্ক্ষিত অন্তিম সিকোয়েন্টটি লিখে শুরু করো।
\begin{prooftree}
\AxiomC{}
\UnaryInf$\lnot !A \lor !B \fCenter !A \lif !B$
\end{prooftree}
এই অন্তিম সিকোয়েন্ট দিতে পারে এমন কোনো যৌক্তিক বিধি খুঁজতে এর যৌক্তিক
সংযোজকগুলি দেখি: $\lnot$, $\lor$ ও~$\lif$। আপাতত শুধু $\lor$ ও~$\lif$
প্রাসঙ্গিক, কারণ তারাই অন্তিম সিকোয়েন্টের !!{sentence}s-এর
!!{main operator}s; $\lnot$ অন্য সংযোজকের পরিধির ভিতরে আছে, তাই সেটির ব্যবস্থা
পরে করব। সুতরাং শেষ অনুমানের যৌক্তিক বিধি হতে পারে \LeftR{\lor} অথবা
\RightR{\lif}। যেকোনোটিই নেওয়া যায়; আমরা \RightR{\lif} নিই—অন্তত
এ কারণে যে এতে আরও কিছুক্ষণ দুই শাখায় ভাগ করা এড়ানো যায়। যে
\RightR{\lif} অনুমানটি এই নিচের সিকোয়েন্ট দেয়, তার রূপ হবে:
\begin{prooftree}
\AxiomC{}
\UnaryInf$ !A, \lnot !A \lor !B \fCenter !B $
\RightLabel{\RightR{\lif}} \UnaryInf$ \lnot !A \lor !B \fCenter !A \lif !B $
\end{prooftree}

$\lnot !A \lor !B$-কে পূর্বাংশের একেবারে বাঁ প্রান্তে আনলে
\LeftR{\lor} বিধি প্রয়োগ করা যায়। বিধি-ছক অনুসারে এটি নিচের মতো দুটি
উপরের সিকোয়েন্টে ভাগ হবে:
\begin{prooftree}
\AxiomC{}
\UnaryInf$\lnot !A, !A \fCenter !B$
\AxiomC{}
\UnaryInf$!B, !A \fCenter !B$
\RightLabel{\LeftR{\lor}}
\BinaryInf$ \lnot !A \lor !B, !A \fCenter !B $
\RightLabel{\LeftR{\Exchange}}
\UnaryInf$ !A, \lnot !A \lor !B \fCenter !B $
\RightLabel{\RightR{\lif}}
\UnaryInf$ \lnot !A \lor !B \fCenter !A \lif !B $
\end{prooftree}
মনে রেখো, আমরা উপর দিকে এগিয়ে প্রারম্ভিক সিকোয়েন্টে পৌঁছতে চাই;
দেখা যাচ্ছে আমরা বেশ কাছাকাছি!{} ডান শাখাটি একটি প্রারম্ভিক সিকোয়েন্ট
থেকে মাত্র একটি দুর্বলীকরণ ও একটি অদলবদল দূরে; তারপরই সেটি সম্পূর্ণ:
\begin{prooftree}
\AxiomC{}
\UnaryInf$\lnot !A, !A \fCenter !B$
\Axiom$!B \fCenter !B$
\RightLabel{\LeftR{\Weakening}}
\UnaryInf$!A, !B \fCenter !B$
\RightLabel{\LeftR{\Exchange}}
\UnaryInf$!B, !A \fCenter !B$
\RightLabel{\LeftR{\lor}}
\BinaryInf$\lnot !A \lor !B, !A \fCenter !B $
\RightLabel{\LeftR{\Exchange}}
\UnaryInf$ !A, \lnot !A \lor !B \fCenter !B $
\RightLabel{\RightR{\lif}}
\UnaryInf$ \lnot !A \lor !B \fCenter !A \lif !B $
\end{prooftree}

এবার বাঁ শাখাটি দেখি। কোনো !!{sentence}-এ একমাত্র যৌক্তিক সংযোজক
হলো পূর্বাংশের !!{sentence}s-এ থাকা~$\lnot$; তাই এখানে \LeftR{\lnot}
বিধির একটি প্রয়োগ চাই।
\begin{prooftree}
\AxiomC{}
\UnaryInf$ !A \fCenter !B, !A$
\RightLabel{\LeftR{\lnot}}
\UnaryInf$\lnot !A, !A \fCenter !B$
\Axiom$!B \fCenter !B$
\RightLabel{\LeftR{\Weakening}}
\UnaryInf$!A, !B \fCenter !B$
\RightLabel{\LeftR{\Exchange}}
\UnaryInf$!B, !A \fCenter !B$
\RightLabel{\LeftR{\lor}}
\BinaryInf$\lnot !A \lor !B, !A \fCenter !B $
\RightLabel{\LeftR{\Exchange}}
\UnaryInf$ !A, \lnot !A \lor !B \fCenter !B $
\RightLabel{\RightR{\lif}}
\UnaryInf$ \lnot !A \lor !B \fCenter !A \lif !B $
\end{prooftree}
ডান শাখাটি যেভাবে শেষ করেছি, একইভাবে মাত্র একটি দুর্বলীকরণ ও একটি
অদলবদল দিয়ে বাঁ শাখাটিও শেষ করা যায়।
\begin{prooftree}
\Axiom$!A \fCenter !A$
\RightLabel{\RightR{\Weakening}}
\UnaryInf$ !A \fCenter !A, !B$
\RightLabel{\RightR{\Exchange}}
\UnaryInf$ !A \fCenter !B, !A$
\RightLabel{\LeftR{\lnot}}
\UnaryInf$\lnot !A, !A \fCenter !B$
\Axiom$!B \fCenter !B$
\RightLabel{\LeftR{\Weakening}}
\UnaryInf$!A, !B \fCenter !B$
\RightLabel{\LeftR{\Exchange}}
\UnaryInf$!B, !A \fCenter !B$
\RightLabel{\LeftR{\lor}}
\BinaryInf$\lnot !A \lor !B, !A \fCenter !B $
\RightLabel{\LeftR{\Exchange}}
\UnaryInf$ !A, \lnot !A \lor !B \fCenter !B $
\RightLabel{\RightR{\lif}}
\UnaryInf$ \lnot !A \lor !B \fCenter !A \lif !B $
\end{prooftree}
% BN-SRC-025 TARGET-CORRECTION-BEGIN
\par\smallskip\noindent\textnormal{\small\textbf{উৎস-সংশোধন (BN-SRC-025):} হিমায়িত ইংরেজি উৎসে উপরের পুনরাবৃত্ত প্রমাণ-বৃক্ষগুলিতে পূর্বাংশের $\lnot !A \lor !B$ ও~$!A$-এর স্থান অদলবদল করা চারটি ধাপকে $\RightR{\Exchange}$ বলা হয়েছে। বদলটি সিকোয়েন্টের বাঁদিকে ঘটেছে; তাই বাংলা বৃক্ষগুলিতে প্রযোজ্য $\LeftR{\Exchange}$ লেবেল ব্যবহার করা হয়েছে।} % BN-SRC-025 TARGET-CORRECTION-END
\end{ex}

\begin{ex}
$\lnot !A \lor \lnot !B \Sequent \lnot (!A \land !B)$ সিকোয়েন্টটির
একটি $\Log{LK}$-!!{derivation} দাও।

আগের কৌশল অনুসারে নিচে কাঙ্ক্ষিত অন্তিম সিকোয়েন্টটি লিখে শুরু করি।
\begin{prooftree}
\AxiomC{}
\UnaryInf$ \lnot !A \lor \lnot !B \fCenter \lnot (!A \land !B) $
\end{prooftree}
অন্তিম সিকোয়েন্টের !!{sentence}s-এর যে মুখ্য সংযোজকগুলি ব্যবহার করা
যায়, সেগুলি হলো $\lor$ ও~$\lnot$। এখানে \LeftR{\lor} বা
\RightR{\lnot}—দুটির যেকোনোটি প্রয়োগ করা যায়। আমরা \RightR{\lnot}
দিয়ে শুরু করি, কারণ তাতে আরও কিছুক্ষণ দুই শাখায় ভাগ করা এড়ানো যায়:
\begin{prooftree}
\AxiomC{}
\UnaryInf$!A \land !B, \lnot !A \lor \lnot !B \fCenter $
\RightLabel{\RightR{\lnot}}
\UnaryInf$\lnot !A \lor \lnot !B \fCenter \lnot (!A \land !B)$
\end{prooftree}
এবার \LeftR{\land} নাকি \LeftR{\lor} বিধি দেখব, সেই পছন্দ আছে।
\LeftR{\land} প্রয়োগ করলে কী হয় দেখা যাক: শুরুতে হয়
$!A, \lnot !A \lor \lnot !B \Sequent \quad$, নয়তো
$!B, \lnot !A \lor \lnot !B \Sequent \quad$ নেওয়া যায়।
!!{derivation}-টি $!A$ ও~$!B$-এর বিনিময়ে প্রতিসম, তাই প্রথমটি নিই:
% BN-SRC-026 TARGET-CORRECTION-BEGIN
\par\smallskip\noindent\textnormal{\small\textbf{উৎস-সংশোধন (BN-SRC-026):} এই দুটি সম্ভাব্য উপরের সিকোয়েন্ট লিখতে গিয়ে হিমায়িত ইংরেজি গদ্যে দ্বিতীয় বিযোজ্যে $\lnot$ বাদ গেছে। উদাহরণের নিয়ামক অন্তিম সিকোয়েন্ট, প্রতিসাম্যের মন্তব্য এবং পরের সব প্রমাণ-বৃক্ষে $\lnot !A \lor \lnot !B$ আছে; বাংলা গদ্যে সেই রূপই পুনরুদ্ধার করা হয়েছে।} % BN-SRC-026 TARGET-CORRECTION-END
\begin{prooftree}
\AxiomC{}
\UnaryInf$!A, \lnot !A \lor \lnot !B \fCenter $
\RightLabel{\LeftR{\land}}
\UnaryInf$!A \land !B, \lnot !A \lor \lnot !B \fCenter $
\RightLabel{\RightR{\lnot}}
\UnaryInf$\lnot !A \lor \lnot !B \fCenter \lnot (!A \land !B)$
\end{prooftree}
!!{derivation}-টি পূরণ করতে থাকলে দেখা যায় একটি সমস্যা হচ্ছে:
\begin{prooftree}
\Axiom$!A \fCenter !A$
\RightLabel{\LeftR{\lnot}}
\UnaryInf$ \lnot !A, !A \fCenter$
\AxiomC{}
\RightLabel{?}
\UnaryInf$!A \fCenter !B$
\RightLabel{\LeftR{\lnot}}
\UnaryInf$ \lnot !B, !A \fCenter$
\RightLabel{\LeftR{\lor}}
\BinaryInf$\lnot !A \lor \lnot !B, !A \fCenter $
\RightLabel{\LeftR{\Exchange}}
\UnaryInf$!A, \lnot !A \lor \lnot !B \fCenter $
\RightLabel{\LeftR{\land}}
\UnaryInf$!A \land !B, \lnot !A \lor \lnot !B \fCenter $
\RightLabel{\RightR{\lnot}}
\UnaryInf$\lnot !A \lor \lnot !B \fCenter \lnot (!A \land !B)$
\end{prooftree}
ডান শাখার শীর্ষ সিকোয়েন্টটি আর সরল করা যায় না এবং গঠনগত অনুমান দিয়েও সেটিকে
কোনো প্রারম্ভিক সিকোয়েন্টে আনা যায় না। তাই এটি সঠিক পথ নয়। স্পষ্টতই
উপরে \LeftR{\land} প্রয়োগ করা ভুল হয়েছিল। আগের অবস্থায় ফিরে তার
বদলে \LeftR{\lor} প্রয়োগ করলে পাই
\begin{prooftree}
\AxiomC{}
\UnaryInf$\lnot !A, !A \land !B \fCenter $

\AxiomC{}
\UnaryInf$\lnot !B, !A \land !B \fCenter $

\RightLabel{\LeftR{\lor}}
\BinaryInf$\lnot !A \lor \lnot !B, !A \land !B \fCenter $
\RightLabel{\LeftR{\Exchange}}
\UnaryInf$!A \land !B, \lnot !A \lor \lnot !B \fCenter $
\RightLabel{\RightR{\lnot}}
\UnaryInf$\lnot !A \lor \lnot !B \fCenter \lnot (!A \land !B)$
\end{prooftree}
প্রতিটি শাখা আগের মতো পূরণ করলে পাই
\begin{prooftree}
\Axiom$ !A \fCenter!A$
\RightLabel{\LeftR{\land}}
\UnaryInf$!A \land !B \fCenter !A$
\RightLabel{\LeftR{\lnot}}
\UnaryInf$\lnot !A, !A \land !B \fCenter $

\Axiom$ !B \fCenter !B$
\RightLabel{\LeftR{\land}}
\UnaryInf$!A \land !B \fCenter !B$
\RightLabel{\LeftR{\lnot}}
\UnaryInf$\lnot !B, !A \land !B \fCenter $

\RightLabel{\LeftR{\lor}}
\BinaryInf$\lnot !A \lor \lnot !B, !A \land !B \fCenter $
\RightLabel{\LeftR{\Exchange}}
\UnaryInf$!A \land !B, \lnot !A \lor \lnot !B \fCenter $
\RightLabel{\RightR{\lnot}}
\UnaryInf$\lnot !A \lor \lnot !B \fCenter \lnot (!A \land !B)$
\end{prooftree}
(এই ধাপগুলিতে $\lnot$ বিধিগুলির নিচে $\land$ বিধিগুলি প্রয়োগ করেও
সঠিক !!{derivation} পাওয়া যেত।)
\end{ex}

\begin{ex}
এ পর্যন্ত সংকোচন-বিধি ব্যবহার করিনি, কিন্তু কখনো এটি দরকার হয়। এমন একটি
উদাহরণ দেখা যাক। ধরা যাক, $\quad \Sequent !A \lor \lnot !A$ প্রমাণ
করতে চাই। নিচ থেকে উপর দিকে $\RightR{\lor}$ প্রয়োগ করলে নিচের দুটি
!!{derivation}s-এর কোনো একটি পাওয়া যাবে:
\begin{prooftree}
\AxiomC{}
\UnaryInf$ \fCenter !A$
\RightLabel{\RightR{\lor}}
\UnaryInf$ \fCenter !A \lor \lnot !A$
\DisplayProof\qquad\bottomAlignProof
\AxiomC{}
\UnaryInf$!A \fCenter $
\RightLabel{\RightR{\lnot}}
\UnaryInf$ \fCenter \lnot !A$
\RightLabel{\RightR{\lor}}
\UnaryInf$ \fCenter !A \lor \lnot !A$
\end{prooftree}
অবশ্য এর কোনোটিই কোনো প্রারম্ভিক সিকোয়েন্টে শেষ হয় না। কৌশলটি হলো
বুঝতে পারা যে সংকোচন-বিধি কোনো !!a{sentence}-এর দুটি প্রতিলিপিকে
একটিতে মেলাতে দেয়—আর যখন নিচ থেকে উপর দিকে প্রমাণ খুঁজি, তখন পূর্বধারণায়
$!A \lor \lnot !A$-এর একটি প্রতিলিপি রেখে দিতে পারি; যেমন,
\begin{prooftree}
\AxiomC{}
\UnaryInf$ \fCenter !A \lor \lnot !A, !A$
\RightLabel{\RightR{\lor}}
\UnaryInf$ \fCenter !A \lor \lnot !A, !A \lor \lnot !A$
\RightLabel{\RightR{\Contraction}}
\UnaryInf$ \fCenter !A \lor \lnot !A$
\end{prooftree}
এবার দ্বিতীয়বার $\RightR{\lor}$ প্রয়োগ করে~$\lnot !A$-ও পাওয়া যায়,
যা একটি সম্পূর্ণ !!{derivation}-এ নিয়ে যায়।
\begin{prooftree}
\Axiom$!A \fCenter !A$
\RightLabel{\RightR{\lnot}}
\UnaryInf$\fCenter !A, \lnot !A$
\RightLabel{\RightR{\lor}}
\UnaryInf$\fCenter !A, !A \lor \lnot !A$
\RightLabel{\RightR{\Exchange}}
\UnaryInf$ \fCenter !A \lor \lnot !A, !A$
\RightLabel{\RightR{\lor}}
\UnaryInf$ \fCenter !A \lor \lnot !A, !A \lor \lnot !A$
\RightLabel{\RightR{\Contraction}}
\UnaryInf$ \fCenter !A \lor \lnot !A$
\end{prooftree}
\end{ex}

\begin{prob}
নিচের সিকোয়েন্টগুলির !!{derivation}s দাও:
\begin{enumerate}
\item $!A \land (!B \land !C) \Sequent (!A \land !B) \land !C$।
\item $!A \lor (!B \lor !C) \Sequent (!A \lor !B) \lor !C$।
\item $!A \lif (!B \lif !C) \Sequent !B \lif (!A \lif !C)$।
\item $!A \Sequent \lnot\lnot !A$।
\end{enumerate}
\end{prob}

\begin{prob}
নিচের সিকোয়েন্টগুলির !!{derivation}s দাও:
\begin{enumerate}
\item $(!A \lor !B) \lif !C \Sequent !A \lif !C$।
\item $(!A \lif !C) \land (!B \lif !C) \Sequent (!A \lor !B) \lif !C$।
\item $\Sequent \lnot(!A \land \lnot !A)$।
\item $!B \lif !A \Sequent \lnot !A \lif \lnot !B$।
\item $\Sequent (!A \lif \lnot !A) \lif \lnot !A$।
\item $\Sequent \lnot(!A \lif !B) \lif \lnot !B$।
\item $!A \lif !C \Sequent \lnot (!A \land \lnot !C)$।
\item $!A \land \lnot !C \Sequent \lnot (!A \lif !C)$।
\item $!A \lor !B, \lnot !B \Sequent !A$।
\item $\lnot !A \lor \lnot !B \Sequent \lnot(!A \land !B)$।
\item $\Sequent (\lnot !A \land \lnot !B) \lif\lnot(!A \lor !B)$।
\item $\Sequent \lnot(!A \lor !B) \lif (\lnot !A \land \lnot !B)$।
\end{enumerate}
\end{prob}

\begin{prob}
নিচের সিকোয়েন্টগুলির !!{derivation}s দাও:
\begin{enumerate}
\item $\lnot(!A \lif !B) \Sequent !A$।
\item $\lnot(!A \land !B) \Sequent \lnot !A \lor \lnot !B$।
\item $!A \lif !B \Sequent \lnot !A \lor !B$।
\item $\Sequent \lnot \lnot !A \lif !A$।
\item $!A \lif !B, \lnot !A \lif !B \Sequent !B$।
\item $(!A \land !B) \lif !C \Sequent (!A \lif !C) \lor (!B \lif !C)$।
\item $(!A \lif !B) \lif !A \Sequent !A$।
\item $\Sequent (!A \lif !B) \lor (!B \lif !C)$।
\end{enumerate}
(এগুলির প্রতিটিতে $\RightR{\Contraction}$~বিধি দরকার।)
\end{prob}
\end{document}
